\chapter{Listado de acrónimos}

\begin{acronym}[XXXXXXXX]
  \acro{API}     {Application Programming Interface}
  \acro{CORS}    {Cross-Origin Resource Sharing}
  \acro{CRUD}    {Create, Read, Update, Delete}
  \acro{CSV}     {Comma-Separated Values}
  \acro{DMBOK}   {Data Management Body of Knowledge}
  \acro{EDA}     {Exploratory Data Analysis}
  \acro{HTTP}    {Hypertext Transfer Protocol}
  \acro{HU}      {Historia de Usuario}
  \acro{IA}      {Inteligencia Artificial}
  \acro{IEC}     {International Electrotechnical Commission}
  \acro{IEEE}    {Institute of Electrical and Electronics Engineers}
  \acro{IQR}     {Interquartile Range}
  \acro{IDOR}    {Insecure Direct Object Reference}
  \acro{ISO}     {International Organization for Standardization}
  \acro{JSON}    {JavaScript Object Notation}
  \acro{JWT}     {JSON Web Token}
  \acro{KPI}     {Key Performance Indicator}
  \acro{ML}      {Machine Learning}
  \acro{MLOps}   {Machine Learning Operations}
  \acro{MVC}     {Model-View-Controller}
  \acro{MVP}     {Minimum Viable Product}
  \acro{ORM}     {Object-Relational Mapping}
  \acro{PBKDF2}  {Password-Based Key Derivation Function 2}
  \acro{PII}     {Personally Identifiable Information}
  \acro{QA}      {Quality Assurance}
  \acro{REST}    {Representational State Transfer}
  \acro{RGPD}    {Reglamento General de Protección de Datos}
  \acro{SaaS}    {Software as a Service}
  \acro{SHA}     {Secure Hash Algorithm}
  \acro{SLA}     {Service Level Agreement}
  \acro{SodaCL}  {Soda Checks Language}
  \acro{SQL}     {Structured Query Language}
  \acro{SQuaRE}  {Systems and Software Quality Requirements and Evaluation}
  \acro{SSO}     {Single Sign-On}
  \acro{TFG}     {Trabajo Fin de Grado}
  \acro{UI}      {User Interface}
  \acro{UML}     {Unified Modeling Language}
  \acro{UX}      {User Experience}
  \acro{WCAG}    {Web Content Accessibility Guidelines}
  \acro{WSGI}    {Web Server Gateway Interface}
  \acro{YAML}    {YAML Ain't Markup Language}
\end{acronym}


% \ac{OO}   la primera vez \acf, después \acs
% \acs{OO}  short: OO
% \acl{OO}  large: Object Oriented
% \acf{OO}  full : Object Oriented (OO)
% \acx{OO}         OO (Object Oriented)

% usa \Acro en lugar de \acro cuando quieras que siempre aparezca en formato 'short'

% Local variables:
%   TeX-master: "main.tex"
% End:
